\documentclass[12pt,a4paper]{report}
\usepackage{fontspec}
\setmainfont{Latin Modern Roman}
\usepackage{arabxetex}
\newfontfamily\arabicfont[Script=Arabic,Scale=1.15]{Traditional Arabic}

\usepackage[a4paper, top=2.5cm, left=2.5cm, right=2.5cm, bottom=3.5cm]{geometry}
\usepackage{setspace}
\onehalfspacing

\begin{document}
\pagenumbering{roman}
\setcounter{page}{3}

\chapter*{%
  \begin{flushright}
    \textarabic{ملخص}
  \end{flushright}%
}
\begin{Arabic}

نادراً ما تنتج الأعطال الكهروضوئية أعراضاً مرئية، ولا يمكن استنتاجها إلا من الانحرافات في الإشارات الكهربائية والأرصاد الجوية المقاسة، مما يجعل الكشف والتشخيص الآلي للأعطال (\LR{FDD}) ضرورياً مع توسع نطاق التركيبات. تطور هذه الأطروحة نظاماً شاملاً للكشف عن الأعطال وتشخيصها قائماً على الذكاء الاصطناعي، تم التحقق من صحته باستخدام بيانات ميدانية حقيقية من محطة \LR{Costa} الكهروضوئية، ونشره على منصة مدمجة.

يعتمد النظام على خط أنابيب بيانات متعدد المراحل (الاستيعاب، التقسيم الزمني، المعالجة المسبقة، وهندسة الميزات) يحول القياسات الخام بتردد \LR{1\,Hz} إلى ميزات مستنيرة بالفيزياء، مع منع صارم لتسرب البيانات: تقسيم زمني مدرك للمقاطع وتقدير الإحصائيات فقط على مجموعة التدريب في كل مرحلة. ثم تُعالج مهمتان متكاملتان. بالنسبة لكشف الشذوذ، يُدرَّب كاشف شبه مراقب على التشغيل الطبيعي فقط، مع سياسة موحدة لتوزيع باريتو المعمم لمعايرة العتبات بشكل متماثل عبر جميع الكاشفات لضمان مقارنة عادلة؛ يحقق النموذج المنشور القائم على التدفق التطبيعي المشروط بالفيزياء (\LR{PC-Flow}، أقل من \LR{25,000} معامل) \LR{PR-AUC} ثنائياً للاختبار قدره \LR{0.989} و\LR{PR-AUC} لأسوأ فئة قدره \LR{0.968} تحت التحقق المتقاطع بخمسة أقسام طبقية حسب النوبات. بالنسبة لتصنيف الأعطال، تُستخدم مجموعات الأشجار المعززة بالتدرج مع ترجيح الفئات، حيث يحقق مصنف \LR{ExtraTrees} المختار \LR{F1} كلياً قدره \LR{0.959} ودقة متوازنة قدرها \LR{\%94.3} تحت التحقق المتقاطع بخمسة أقسام. تُصدَّر النماذج المُتحقق منها إلى صيغة \LR{ONNX} وتُشغَّل على منصة \LR{NVIDIA Jetson Nano} خلف واجهة مراقبة في الوقت الفعلي، حيث يحول نظام الإنذار الهرمي ثلاثي المستويات وفق معيار \LR{ISA-18.2} (\LR{P3} حساس، \LR{P2} مطابق، \LR{P1 CUSUM}) درجة الكاشف إلى إنذارات ذات أولوية، إلى جانب تشخيص لكل فئة ومراقبة الانحراف بطريقة \LR{Page-Hinkley} مع تحديثات للنموذج بموافقة المشغل.

ترسم هذه النتائج مجتمعة مساراً قابلاً للتكرار من القياسات الكهروضوئية الواقعية إلى وعي موثوق بالأعطال على الجهاز المدمج.

\vspace{1cm}

\noindent\rule[2pt]{\textwidth}{0.5pt}
\textbf{الكلمات المفتاحية:} الأنظمة الكهروضوئية، كشف وتشخيص الأعطال، التعلم الآلي، التعلم العميق، الصيانة التنبؤية.

\noindent\rule[2pt]{\textwidth}{0.5pt}

\end{Arabic}
\end{document}
